\documentclass[9pt,twoside]{extarticle}
\usepackage{Breviary-Appendix}
\usepackage{indentfirst}
%%
\begin{document}
\renewcommand{\thepage}{\arabic{page}*}
\begin{center} \includegraphics[width=1 \textwidth]{images/Sacred_Hearts_rotated} \end{center}
\subsection*{\center \normalfont \textsc{die 22 augusti}}

\medskip
\section*{\center \cloisterblack \textred{I}mmaculati \textred{C}ordis \textred{B}eatæ \textred{M}ariæ \textred{V}irginis}
%\begin{center}\textred{Virginis Ordinis Nostri.}\end{center}
\subsubsection*{\center \boldred{DUPLEX MAJUS II CLASSIS}}
\fancyhead[CO]{\textsc{\textred{Immaculati Cordis B. Mariæ Virginis}}}
\fancyhead[CE]{\textsc{\textred{Die 22 Augusti}}}

\begin{multicols}{2}

{\small \textred{Omnia de Communi Festorum B. Mariæ Virginis [161], præter ea quæ hic habentur propria.}}

\subsubsection*{\center \normalfont \textsc{IN I VESPERIS}}

\subsubsection*{\normalfont \textred{Capitulum \hfill Cant.~8}}
\lettrineThree{P}{o}ne me ut signáculum super cor tuum, ut signáculum super bráchium tuum, quia fortis est ut mors diléctio, dura sicut inférnus æmulátio. Lámpades ejus lámpades ignis, atque flammárum.

\smallskip
\textred{℟.}~Ego Mater pulchræ dilectiónis, et timóris, et agnitiónis et sanctæ spei; \textred{*} In me grátia omnis viæ, et veritátis. \textred{℣.}~In me omnis spes vitæ et virtútis. \textred{I}n me. Glória Patri. \textred{I}n me.

\textred{Hymnus.} Ave maris stella. \textred{[164].}

\textred{℣.}~Exsúlta in omni corde fília Jerúsalem. \textred{℟.}~Rex Israel Dóminus in médio tui, allelúja.

\textred{Ad Magnif. Ant.} Exsultávit Cor meum * in Dómino, et exaltátum est cornu meum in Deo meo~: quia lætáta sum in salutári tuo.

\subsubsection*{\center \normalfont \textred{Oratio}}
\lettrine{O}{m}nípotens sempitérne Deus, qui in Corde beátæ Maríæ Vírginis dignum Spíritus sancti habitáculum præparásti, concéde propítius, ut ejúsdem immaculáti Cordis festivitátem devóte mente recoléntes, secúndum Cor tuum vívere valeámus. Per Dóminum...in unitáte ejúsdem Spíritus.

%\small \textred{Et fit commem. Domin. occurrentis.}
\small
\textred{Et fit Commemoratio præcedentis~:} \normalsize

\textred{Ant.} Símile est * regnum cælórum hómini negotiatóri quærénti bonas margarítas~: invénta una pretiósa margaríta, dedit ómnia sua, et comparávit eam.

\textred{℣.}~Diffúsa est grátia in lábiis tuis. \textred{℟.}~Proptérea benedíxit te Deus in ætérnum, allelúja.

\noindent Orémus. \hfill \textred{Oratio}
\lettrine{O}mnípotens et miséricors Deus, qui beátam Joánnam Francíscam tuo amóre succénsam, admirábili spíritus fortitúdine per omnes vitæ sémitas in via perfectiónis donásti, quique per illam illustráre Ecclésiam tuam nova prole voluísti: ejus méritis et précibus concéde; ut, qui infirmitátis nostræ cónscii de tua virtúte confídimus, cæléstis grátiæ auxílio cuncta nobis adversántia vincámus. Per Dóminum.

\small \textred{De Octavo Assumptionis BMV nihil fit.}

\textred{Completorium de Dominica.}

\textred{In fine Hymnorum Horarum et Completorii dicitur ℣.} Glória tibi Dómine, Qui natus es de Vírgine.

\textred{Ad Nunc dimittis, Ant.} Allelúja~: Sancta Dei Génitrix.
\normalsize

\subsubsection*{\center \normalfont \textsc{AD MATUTINUM}}
\indent \textred{\small Sequens Invitatorium bis ante Psalmum et in repetitione post Psalmum dicitur flexis genibus.}

\textred{Invit.} Ave María, grátia plena, \textred{*} Dóminus tecum.

\textred{Ps. 94} Veníte, exsultémus. \textred{2.}

\textred{Hymn.} Quem terra, \textred{de Communi [165].}

\subsubsection*{\center \normalfont \small \textsc{IN I NOCTURNO}}

\smallskip
\textred{Ant.} Benedícta tu * inter mulíeres, et benedíctus fructus ventris tui.

%{\small \textred{Ps.} Dómine Dóminus noster, \textred{cum reliquis de commune B.M.V. [165].}}
{\small \textred{Psalmi de Communi [165].}}

\textred{Ant.} Et ait María~: * Magníficat ánima mea Dóminum, et exsultávit spíritus meus in Deo salutári meo.

\textred{Ant.} Quia respéxit * humilitátem ancíllæ suæ; ecce enim ex hoc beátam me dicent omnes generatiónes.

\smallskip
\textred{℣.}~Exsúlta in omni corde, fília Jerúsalem. \textred{℟.}~Rex Israel Dóminus in médio tui, allelúja.

\smallskip
{\small \textred{Bened.} Alma Virgo vírginum intercédat pro nobis ad Dóminum.}

\subsubsection*{\center \normalfont De Cánticis Canticórum}
\subsubsection*{\center \normalfont \textred{Lectio j \hfill Cap.~4,~1-4,~9,~11-12}}
\lettrineThree{Q}{u}am pulchra es amíca mea, quam pulchra es! Oculi tui columbárum, absque eo quod intrínsecus latet. 
Sicut vitta coccínea lábia tua, et elóquium tuum dulce. Sicut fragmen mali púnici, ita genæ tuæ, absque eo quod intrínsecus latet.
Sicut turris David collum tuum, quæ ædificáta est cum propugnáculis~: mille clýpei pendent ex ea, omnis armatúra fórtium.
Vulnerásti cor meum, soror mea sponsa, vulnerásti cor meum in uno oculórum tuórum, et in uno crine colli tui.
Favus distíllans lábia tua Sponsa~: mel et lac sub lingua tua.
Hortus conclúsus, fons signátus.

\textred{℟.}~Deduc me Dómine, in via tua, et ingrédiar in veritáte tua; lætétur cor meum, \textred{*} Ut tímeat nomen tuum. \textred{℣.}~Pone me ut signáculum super cor tuum, ut signáculum super bráchium tuum. \textred{U}t tímeat.

\smallskip
{\small \textred{Bened.} Nos cum prole pia benedícat Virgo María.}

\subsubsection*{\center \normalfont \textred{Lectio ij \hfill Cap.~5,~1-6}}
\lettrine{V}{e}ni in hortum meum soror mea sponsa~: méssui myrrham meam cum aromátibus meis~: comédi favum cum melle meo.
Vox dilécti mei pulsántis~: Aperi mihi, soror mea, amíca mea, colúmba mea, immaculáta mea~: quia caput meum plenum est rore, et cincínni mei guttis nóctium.
Surréxi ut aperírem dilécto meo~: manus meæ stillavérunt myrrham, et dígiti mei pleni myrrha probatíssima.
Apérui dilécto meo~: at ille declináverat, atque transíerat. Anima mea liquefácta est, ut locútus est.

\textred{℟.}~In Dómino gaudébo, et exsultábo in Deo Jesu meo~: Deus Dóminus fortitúdo mea, \textred{*} Super excélsa dedúcet me. \textred{℣.}~Læva ejus sub cápite meo, et déxtera illíus amplexábitur me. \textred{S}uper.

\smallskip
{\small \textred{Bened.} Sancta Dei Génitrix sit nobis auxiliátrix.}

\subsubsection*{\center \normalfont \textred{Lectio iij \hfill Cap.~5,~9-16~et~6,~1-9}}
\lettrine{Q}{u}alis est diléctus tuus ex dilécto, o pulchérrima mulíerum? qualis est diléctus tuus ex dilécto, quia sic adjurásti nos?
Diléctus meus cándidus et rubicúndus, eléctus ex míllibus, totus desiderábilis. Talis est diléctus meus, et ipse est amícus meus, fíliæ Jerúsalem.
Quo ábiit diléctus tuus, o pulchérrima mulíerum? quo declinávit diléctus tuus? et quærémus eum tecum.
Pulchra es amíca mea, suávis, et decóra sicut Jerúsalem, terríbilis ut castrórum ácies ordináta.
Una est colúmba mea, perfécta mea, una est matris suæ, elécta genetríci suæ. Vidérunt eam fíliæ, et beatíssimam prædicavérunt~: regínæ laudavérunt eam.

\textred{℟.}~Invéni grátiam apud óculos tuos, Dómine mi, \textred{*} Qui consolátus es me, et locútus es ad Cor ancíllæ tuæ. \textred{℣.}~In toto corde meo exquisívi te, ne repéllas me a mandátis tuis. \textred{Q}ui. Glória Patri. \textred{Q}ui.

\subsubsection*{\center \normalfont \small \textsc{IN II NOCTURNO}}

\smallskip
\textred{Ant.} Beáta, * quæ credidísti, quóniam perficiéntur ea, quæ dicta sunt tibi a Dómino.

\textred{Ant.} Fecit mihi magna * qui potens est, et sanctum nomen ejus.

\textred{Ant.} Fecit poténtiam * in bráchio suo, dispérsit supérbos mente cordis sui.

\smallskip
\textred{℣.}~Anima mea liquefácta est. \textred{℟.}~Ut locútus est diléctus, allelúja.

\smallskip
{\small \textred{Ben.} Précibus suæ Matris benedícat nos Fílius Dei Patris.}

\subsubsection*{\center \normalfont Sermo sancti Bernardíni Senénsis\\ \textred{{\small Ex serm. 9. de Visitat.}\\ Lectio iv}}
\lettrineThree{Q}{u}is mortálium, nisi divíno fretus oráculo, de vera Dei et hóminis Genitríce quidquam módicum, sive grande præsúmat pollútis lábiis nomináre~: quam Pater ante sæcula Deus perpétuam prædestinávit in Vírginem digníssimam; Fílius elégit in Matrem; Spíritus sanctus omnis grátiæ domicílium præparávit? Quibus verbis ego homúnculus sensus altíssimos Virgínei Cordis, sanctíssimo ore prolátos éfferam; quibus non súfficit lingua ómnium hóminum et Angelórum? Dóminus enim ait~: Bonus homo de bono thesáuro cordis sui profert bona. Quis autem inter puros hómines mélior homo potest excogitári, quam illa, quæ méruit éffici Mater Dei, quæ in Corde et in útero suo ipsum Deum hospitáta est. Quis thesáurus mélior, quam ipse divínus amor, quo Cor Vírginis ardens erat?

\textred{℟.}~Ego diligéntes me díligo, et qui mane vígilant ad me, invénient me. Mecum sunt divítiæ et glória~: \textred{*} Opes supérbæ et justítia. \textred{℣.}~Ut ditem diligéntes me, et thesáuros eórum répleam. \textred{O}pes.

\smallskip
{\small \textred{Ben.} Jesus Maríæ Fílius sit nobis clemens et propítius.}

\subsubsection*{\center \normalfont \textred{Lectio v}}
\lettrine{D}{e} hoc ígitur Corde quasi fornáce divíni ardóris Virgo beáta prótulit verba bona, id est verba ardentíssimæ caritátis. Sicut enim a vase óptimo vino pleno, non potest exíre nisi óptimum vinum~: aut sicut a fornáce summi ardóris non egréditur nisi incéndium fervens~: sic de Corde Matris Christi exíre non pótuit verbum, nisi summi, summéque divíni amóris, atque ardóris. Septem verba tantum miræ senténtiæ et virtútis a Christi benedictíssima Matre legúntur dicta. Cum Angelo bis tantum locúta est. Cum Elísabeth, bis étiam. Cum Fílio étiam bis. Cum minístris in nútiis semel. Hæc septem verba, secúndum septem amóris actus, sub miro gradu et órdine proláta, quasi sunt septem flammæ Cordis ejus.

\smallskip
\textred{℟.}~Transíte ad me omnes, qui concupíscitis me, et a generatiónibus meis implémini. \textred{*} In me grátia omnis viæ et veritátis. \textred{℣.}~Qui me invénerit invéniet vitam, et háuriet salútem a Dómino. \textred{I}n me.

\smallskip
{\small \textred{Bened.} Fílius Vírginis Maríæ det nobis gáudia vitæ.}

\subsubsection*{\center \normalfont \textred{Lectio vi} \hfill Ex ecclesiásticis documéntis}
\lettrine{C}ultum litúrgicum, quo Cordi Immaculáto Vírginis Maríæ débitus tribúitur honor, cuíque plures viri sancti ac mulíeres viam parárunt, ipsa Apostólica Sedes primum approbávit ineúnte sǽculo undevicésimo, cum Pius Papa séptimus festum Puríssimi Cordis Maríæ Vírginis instítuit, ab ómnibus diœcésibus et religiósis famíliis, quæ id petiíssent, pie sanctéque agéndum: quod póstmodum Pius Papa nonus Offício ac Missa própria auxit. Ardens autem stúdium atque optátum, jam sǽculo décimo séptimo exórtum et in dies invaléscens, ut nempe ejúsmodi festum, majóri solemnitáte donátum, totíus Ecclésiæ commúne efficerétur, Summus Póntifex Pius duodécimus benígne excípiens anno millésimo nongentésimo quadragésimo secúndo, bello atrocíssimo per orbem fere totum ingravescénte, infínitas populórum ærúmnas míserans, pro sua in Matrem cæléstem pietáte ac fidúcia genus hóminum univérsum illíus Cordi benigníssimo obsecratióne solémni eníxe commendávit, atque in honórem ejúsdem Immaculáti Cordis festum cum Offício et Missa própriis in perpétuum ubíque celebrándum indíxit.
%\subsubsection*{\center \normalfont \textred{Lectio vj}}
%\lettrine{Q}{u}æ verba consíderans mens devóta, cum ipso Prophéta ait~: Quam dúlcia fáucibus meis, id est ómnibus affectiónibus, elóquia tua! Porro dulcédo hæc quam in verbis beátæ Vírginis devóta mens sentit, est ardor piæ devotiónis quam in eis ánima experimentáliter probat. Distinguámus ígitur per órdinem has septem flammas amóris verbórum Vírginis benedíctæ. Prima est flamma amóris separántis. Secúnda amóris transformántis. Tértia amóris communicántis. Quarta amóris jubilántis. Quinta amóris soporántis. Sexta amóris compatiéntis. Séptima amóris consummántis.

\smallskip
\textred{℟.}~Beátus homo, qui audit me, et qui vígilat ad fores meas quotídie et obsérvat ad postes óstii mei. \textred{*} Qui me invénerit invéniet vitam, et háuriet salútem a Dómino. \textred{℣.}~Qui autem in me peccáverit lædet ánimam suam; omnes qui me odérunt, díligunt mortem. \textred{Q}ui me. Glória Patri. \textred{Q}ui me.

\subsubsection*{\center \small \textsc{IN III NOCTURNO}}

\smallskip
\textred{Ant.} Unde hoc mihi, * ut véniat Mater Dómini mei ad me?

\textred{Ant.} Misericórdia ejus * a progénie in progénies timéntibus eum.

\textred{Ant.} Suscépit Israel * púerum suum, recordátus misericórdiæ suæ.

\smallskip
\textred{℣.}~Ego dilécto meo. \textred{℟.}~Et ad me convérsio ejus, allelúja.

\smallskip
{\small \textred{Bened.} Fons Evangélii répleat nos dógmate cæli.}


\subsubsection*{\center \normalfont Lectio sancti Evangélii\\ secúndum Lucam}
\subsubsection*{\center \normalfont \textred{Lectio vij \hfill Cap.~2,~48-51}}
\lettrine{I}{n} illo témpore~: Dixit Mater Jesu ad illum~: Fili, quid fecísti nobis sic? Ecce, pater tuus et ego doléntes quærebámus te. Et réliqua.

\subsubsection*{\center \normalfont Homilía sancti Bedæ\\ Venerábilis Presbyteri\\ \small \textred{In Homil. de Tempore. Dom. I\\ post Epiph.}}

\smallskip
\lettrineThree{M}{a}ter ejus, inquit Evangelísta, conservábat ómnia verba hæc in Corde suo. Omnia quæ de Dómino vel a Dómino dicta, sive acta cognóverat Mater Virgo, diligéntius in Corde retinébat, sollícite cuncta memóriæ commendábat, ut cum demum tempus prædicándæ sive scribéndæ Incarnatiónis ejus adveníret, sufficiénter univérsa, prout essent gesta, posset explicáre quæréntibus. Imitémur et nos, fratres mei, piam Dómini Matrem, ipsi quoque ómnia verba et facta Dómini ac Salvatóris nostri fixa in corde conservándo.

\smallskip
\textred{℟.}~Aquæ multæ non potuérunt exstínguere caritátem~: \textred{*} Neque flúmina óbruent eam. \textred{℣.}~Si déderit homo omnem substántiam domus suæ pro dilectióne, quasi nihil despíciet eam. \textred{N}eque flúmina.

\smallskip
\small \textred{Bened.} Ab hoste malígno deféndat nos Dei Génitrix Virgo.

%\textred{Si Hom. Dom. non legitur, Lectio octava dividatur in octavam et nonam, et ante Lect. ix. dicitur Ben.} In omni tribulatióne et angústia succúrrat nobis pia Virgo María.
\normalsize
\subsubsection*{\center \normalfont \textred{Lectio viij}}
\lettrine{H}{o}rum Dómini verbórum et actórum meditatióne diúrna et noctúrna importúnos inánium nocentiúmque cogitatiónum repellámus incúrsus; horum crebra collatióne et nos et próximos nostros a fábulis supervácuis, et male dulcorátis detractiónum collóquiis castigáre, atque ad divínæ frequéntiam laudis accéndere curémus.
Si enim, fratres caríssimi, in futúri beatitúdine sæculi habitáre in domo Dómini, ac perpétuo illum laudáre desiderámus, opórtet nimírum ut in hoc sæculo, quod in futúro quærámus, sollícite præmonstrémus~: frequentándo vidélicet Ecclésiæ límina, et non solum in ea laudes Dómini canéndo, sed et in omni loco dominatiónis ejus ea, quæ ad laudem gloriámque Conditóris nostri profíciant, verbis páriter ac factis ostendéndo.

\smallskip
\textred{℟.}~Spíritus meus super mel dulcis et heréditas mea super mel et favum. Qui audit me non confundétur~: \textred{*} Et qui operántur in me non peccábunt. \textred{℣.}~Qui elúcidant me vitam ætérnam habébunt. \textred{E}t qui.

\smallskip
\small \textred{Ben.} In omni tribulatióne et angústia succúrrat nobis pia Virgo María.

\textred{Si hodie fuerit Sabbatum, Lectio ix de Homilia Vigiliæ anticipatæ S. Bartholomæi Apostoli, et ejus Commemoratio ad Laudes, ut die sequenti notatur, 962; et ante Lectionem dicitur Bened.} Verba Redemptóris pia sint suffrágia nobis. \normalsize

\subsubsection*{\center \normalfont Pro sanctis Timótheo et\\ Sóciis Martýribus\\
\textred{Lectio ix}}
\lettrine{T}imótheus Antiochénus, Romam véniens Melchíade summo Pontífice, cum per annum ibi Christi fidem prædicásset, a Tarquínio Urbis præfécto conjícitur in víncula; et, post diutúrnas cárceris ærúmnas, ad idóla perdúcitur, ut eis sacríficet. Quam impietátem summa libertáte detestátus, acerbíssime cǽditur, et excarnificátum corpus viva calce perfúnditur. In quibus aliísque supplíciis constans Martyr, cápite pléctitur. Corpus via Ostiénsi proper sepúlcrum beáti Pauli Apóstoli sepelítur. Quo die étiam, Alexándro imperatóre, apud Ostia Tiberína Hippólytus, epíscopus Portuénsis, ob præcláram fídei confessiónem mánibus pedibúsque ligátis, in altam fóveam aquis plenam præcipitátus, martýrio coronátus est, et ibídem a Christiánis sepúltus. Quo item die, Aureliáno imperatóre, propter eámdem fidem, Augustodúni, Symphoriánus adoléscens várie tortus est. Qui, dum ad últimum supplícium ducerétur, matrem ita clamántem áudiens~: Nate, nate, meménto ætérnæ vitæ, cælum súspice, et ibi regnántem intuére; tibi enim vita non erípitur, sed mutátur in mélius; fórtiter Jesu Christi causa carnífici collum prǽbuit.

%\subsubsection*{\center \textred{Lectio ix}}
%\lettrine{S}{i} enim, fratres caríssimi, in futúri beatitúdine sæculi habitáre in domo Dómini, ac perpétuo illum laudáre desiderámus, opórtet nimírum ut in hoc sæculo, quod in futúro quærámus, sollícite præmonstrémus~: frequentándo vidélicet Ecclésiæ límina, et non solum in ea laudes Dómini canéndo, sed et in omni loco dominatiónis ejus ea, quæ ad laudem gloriámque Conditóris nostri profíciant, verbis páriter ac factis ostendéndo.

%\smallskip
%{\small \textred{Si ix Lect. legitur de Hom. Dom. occurr., Bened.} Verba Redemptóris pia sint suffrágia nobis.}

\textred{℟.}~Ego Mater pulchræ dilectiónis et timóris et agnitiónis et sanctæ spei~: \textred{*} In me grátia omnis viæ et veritátis. \textred{℣.}~In me omnis spes vitæ et virtútis. \textred{I}n me. Glória Patri. \textred{I}n me.

Te Deum laudámus. \textred{6.}

\textred{Sacerd. ℣.}~Viam mandatórum tuórum cucúrri. \textred{℟.}~Cum dilatásti Cor meum, allelúja.

\subsubsection*{\center \normalfont \textsc{AD LAUDES}\\ \normalfont \textred{et per Horas}}

\smallskip
\indent \textred{Ant. 1} Trahe me * post te, currémus in odórem unguentórum tuórum.

\textred{Ps. de Dominica. 27.}

\textred{2} Diléctus meus * lóquitur mihi~: surge, própera amíca mea.

\textred{3} Fulcíte me flóribus * stipáte me malis~: quia amóre lángueo.

\textred{4} Adjúro vos, * fíliæ Jerúsalem, ne suscitétis diléctam quoadúsque ipsa velit.

\textred{5} Ego dórmio, * et cor meum vígilat.


\subsubsection*{\normalfont \textred{Capitulum \hfill Cant.~8}}
\lettrine{P}{o}ne me ut signáculum super cor tuum, ut signáculum super bráchium tuum, quia fortis est ut mors diléctio, dura sicut inférnus æmulátio. Lámpades ejus lámpades ignis, atque flammárum.

\smallskip
\textred{Hymn.} O gloriósa, \textred{de communi B.M.V. [176].}

\textred{℣.}~Qui me invénerit, invéniet vitam. \textred{℟.}~Et háuriet salútem a Dómino, allelúja.

\textred{Ad Bened. Ant.} Extóllens vocem~* quædam múlier de turba, dixit illi~: Beátus venter qui te portávit, et úbera quæ suxísti.

\subsubsection*{\center \normalfont \textred{Oratio}}
\lettrine{O}{m}nípotens sempitérne Deus, qui in Corde beátæ Maríæ Vírginis dignum Spíritus sancti habitáculum præparásti, concéde propítius, ut ejúsdem immaculáti Cordis festivitátem devóte mente recoléntes, secúndum Cor tuum vívere valeámus. Per Dóminum...in unitáte ejúsdem Spíritus.

\smallskip
{\small \textred{Fit Commemoratio Ss. Timothei, Hippolyti et Symphoriani Mm.~:}}

\textred{Ant.} Istórum est enim * regnum cælórum, qui contempsérunt vitam mundi, et pervenérunt ad prǽmia regni, et lavérunt stolas suas in sánguine Agni.

\smallskip
\textred{℣.}~Mirábilis Deus in Sanctis suis. \textred{℟.}~Et gloriósus in majestáte sua.

\smallskip
\noindent Orémus. \hfill \textred{Oratio}
\lettrine{A}uxílium tuum nobis, Dómine, quǽsumus, placátus impénde~: et, intercedéntibus beátis Martýribus tuis Timótheo, Hippóyto et Symphoriáno, déxteram super nos tuæ propitiatiónis exténde. Per Dóminum.

%{\small \textred{Et fit comm. Domin. occurrentis.}}

\subsubsection*{\center \normalfont \textsc{AD PRIMAM}}
\textred{Ant.} Trahe me * post te, currémus in odórem unguentórum tuórum.

{\small \textred{In Responsorio ℣.}~Qui de Vírgine dignátus es nasci.}

\subsubsection*{\center \normalfont \textsc{AD TERTIAM}}
\textred{Ant.} Diléctus meus * lóquitur mihi~: surge, própera amíca mea.

\smallskip
\subsubsection*{\normalfont \textred{Capitulum \hfill Cant.~8}}
\lettrine{P}{o}ne me ut signáculum super cor tuum, ut signáculum super bráchium tuum, quia fortis est ut mors diléctio, dura sicut inférnus æmulátio. Lámpades ejus lámpades ignis, atque flammárum.

\textred{℟.}~Exsúlta in omni corde fília Jerúsalem, \textred{*} Allelúja, allelúja. \textred{E}xsúlta.
\textred{℣.}~Rex Israel Dóminus in médio tui. \textred{A}llelúja, allelúja.
Glória Patri. \textred{E}xsúlta.

\textred{℣.}~Anima mea liquefácta est. \textred{℟.}~Ut locútus est diléctus, allelúja.

\subsubsection*{\center \normalfont \textsc{AD SEXTAM}}
\textred{Ant.} Fulcíte me flóribus * stipáte me malis~: quia amóre lángueo.

\smallskip
\subsubsection*{\normalfont \textred{Capitulum \hfill Eccl.~24}}
\lettrine{E}{go} ex ore Altíssimi prodívi, primogénita ante omnem creatúram, ego feci in cœlis, ut orirétur lumen indefíciens, et sicut nébula texi omnem terram.

\item \textred{℟.}~Anima mea liquefácta est, \textred{*} Allelúja, allelúja. \textred{A}nima mea.
\textred{℣.}~Ut locútus est diléctus. \textred{A}llelúja, allelúja.
Glória Patri. \textred{A}nima mea.

\textred{℣.}~Ego dilécto meo. \textred{℟.}~Et ad me convérsio ejus, allelúja.

\subsubsection*{\center \normalfont \textsc{AD NONAM}}
\textred{Ant.} Ego dórmio, * et cor meum vígilat.

\smallskip
\subsubsection*{\normalfont \textred{Capitulum \hfill Eccl.~24}}
\lettrine{I}{n} omni terra steti, et in omni pópulo, et in ómnibus réquiem quæsívi, et in hereditáte Dómini morábor.

\textred{℟.}~Ego dilécto meo, \textred{*} Allelúja, allelúja. \textred{E}go.
\textred{℣.}~Et ad me convérsio ejus. \textred{A}llelúja, allelúja.
Glória Patri. \textred{E}go.

\textred{℣.}~Qui me invénerit, invéniet vitam. \textred{℟.}~Et háuriet salútem a Dómino, allelúja.

\subsubsection*{\center \normalfont \textsc{IN II VESPERIS}}
{\small \textred{Antiphonæ Laudum, Psalmi de Comm. B.M.V. [177].}}

\subsubsection*{\normalfont \textred{Capitulum \hfill Cant.~8}}
\lettrine{P}{o}ne me ut signáculum super cor tuum, ut signáculum super bráchium tuum, quia fortis est ut mors diléctio, dura sicut inférnus æmulátio. Lámpades ejus lámpades ignis, atque flammárum.

\textred{Hymnus.} Ave maris stella. \textred{[164].}

\textred{℣.}~Exsúlta in omni corde fília Jerúsalem. \textred{℟.}~Rex Israel Dóminus in médio tui, allelúja.

\normalsize \textred{Ad Magnif. Ant.} Exsultávit Cor meum * in Deo salutári meo~: quia fecit mihi magna qui potens est.

\subsubsection*{\center \normalfont \textred{Oratio}}
\lettrine{O}{m}nípotens sempitérne Deus, qui in Corde beátæ Maríæ Vírginis dignum Spíritus sancti habitáculum præparásti, concéde propítius, ut ejúsdem immaculáti Cordis festivitátem devóte mente recoléntes, secúndum Cor tuum vívere valeámus. Per Dóminum...in unitáte ejúsdem Spíritus.
%\small \textred{Et fit comm. Dom. occurr.}

\smallskip 
\small \textred{Et fit Commemoratio sequentis. 960.}
\medskip
\begin{center}\textcolor{red}{\rule{\linewidth}{0.5pt}} \end{center}

\end{multicols}

\end{document}